\documentclass[lualatex,ja=standard]{bxjsarticle}

% --- パッケージ ---
\usepackage{amsmath, amssymb} % 数式用
\usepackage{physics}          % 物理用（\dd や \pdv が便利）
\usepackage{graphicx}         % 画像用
\usepackage{siunitx}          % 単位用（\qty{9.8}{m/s^2} など）
\usepackage{booktabs}         % 綺麗な表用

% --- タイトル情報 ---
\title{実験題目：〇〇の測定}
\author{氏名：〇〇 〇〇（学籍番号：XXXXXXX）}
\date{\today}

\begin{document}
\maketitle

\section{目的}
本実験の目的は……

\section{原理}
ニュートンの運動方程式は以下のように表される。
\begin{equation}
    m \frac{\dd^2 \vb{r}}{\dd t^2} = \vb{F}
\end{equation}

\section{実験方法}
図\ref{fig:setup}に示す装置を用いて……

\section{結果・考察}
実験データを表\ref{table:data}に示す。

\section{結論}
以上の結果より……

\end{document}